\subsection{Background and Motivation}
\label{subsec:background}

The dominant paradigm in digital advertising has, for two decades, relied on a
fundamental assumption that the user interacts with a visual interface (ie web page, 
search results listing, a social media feed) where rectangular display regions
can be allocated to sponsored content. Banner advertisements, programmatic
display auctions, and search-engine sponsored links all exploit this property.
The advertiser purchases screen real estate, the publisher then renders a creative
unit and the user either attends to it or does not. This model generates
hundreds of billions of dollars in annual revenue worldwide and underpins the
economic viability of most free-to-use internet services.

Conversational AI interfaces disrupt this assumption at its foundation. Systems
such as ChatGPT, Claude, and Gemini present information as a single stream of
natural-language text. There are no sidebars, no interstitial slots, and no
visual regions in which a banner can be placed without breaking the
conversational contract between the system and its user. At the same time, these
interfaces are absorbing an increasing share of the queries that users
previously directed to search engines. They include product research, comparison shopping,
recommendation requests, precisely the query types that carry the highest
commercial intent and, consequently, the highest advertising value. The shift is
not merely a change of medium. It is a structural transformation in how users
express intent, receive information, and make purchasing decisions.

The advertising industry therefore faces a concrete technical problem: how to
monetise conversational AI interactions without degrading the user experience
that makes them valuable in the first place. Appending a block of sponsored
links to the end of a chat response replicates the search-ad format but
sacrifices the natural feel of the conversation. Inserting interstitial
advertisements between turns interrupts the dialogue flow and erodes user trust.
Neither approach exploits the distinctive capability of a large language model
(LLM), namely its ability to generate contextually appropriate, fluent
natural-language text that integrates information from multiple sources, including
sponsored product information into a single coherent response.

This observation motivates the present work. If the LLM can understand the
user's intent, match it against an inventory of relevant products, and weave a
sponsored recommendation naturally into an otherwise helpful response, then the
advertisement ceases to be an interruption and becomes, instead, part of the
answer. The result is an advertising format that is native to the conversational
medium: contextually relevant, linguistically integrated, and measurable not
merely through impressions and clicks but through the user's subsequent
conversational engagement with the recommended product.

The same logic extends beyond text. Conversational AI interfaces increasingly
support image generation as a native modality, and a user who asks the system
to render a scene is expressing creative and often commercial intent in a form
that is no less monetisable than a textual product query. If the language model
can weave a sponsored product into a written answer, an image model can, by
analogous reasoning, place that product naturally within a generated scene rather
than overlaying a banner on top of the artwork. Treating text and image
generation as two surfaces of a single advertising pipeline is therefore a
natural extension of the core thesis, and one that this project pursues
explicitly.

A workable platform, however, cannot consist of the generative pipeline alone.
A production-grade advertising system also requires the operational scaffolding
that makes it usable by real advertisers: a way to onboard them, a way to
organise their spend into campaigns with start and end dates and daily
pacing, a way to create and manage ads without touching the database directly,
and a way for both the platform operator and the advertiser to observe how
their money is being spent. These concerns are conventionally relegated to
``future work,'' but they fundamentally shape the data model and the request
path. Budget gating, for example, must happen inside the ad-retrieval loop,
not as an afterthought. This project therefore treats them as first-class
deliverables rather than deferred extensions.

\subsection{Project Objectives}
\label{subsec:objectives}

This project designs, implements, and evaluates an end-to-end LLM-powered
conversational advertising platform (referred to throughout this report as
\textit{Ad LLM as a Service})that delivers inline sponsored product
recommendations within AI chat responses. The platform satisfies the following
objectives:

\begin{enumerate}[label=(\roman*), leftmargin=2.5em]
  \item \textbf{Intent-aware context extraction.} A lightweight LLM call
    extracts structured context from each user message and the preceding
    conversation history, producing a classification of intent, topical
    relevance, sentiment, purchase signal strength, and a
    sensitivity-aware ad receptivity score that gates whether any advertisement
    should be shown at all.

  \item \textbf{Semantic ad matching with budget-aware eligibility.} The
    extracted context is embedded into a vector representation and matched
    against a pre-embedded ad inventory via cosine similarity search over
    \texttt{pgvector}. Candidate ads are then filtered through a campaign
    eligibility check. This verifies that the parent campaign is active, within
    its scheduled window, and has not exhausted either its total budget or
    its daily cap and re-ranked by a business-rule stage that incorporates
    bid price, per-session frequency caps, and brand-safety tags.

  \item \textbf{Natural ad integration with mandatory disclosure.} The main
    response-generation LLM receives the matched advertisement as part of its
    prompt context and is instructed to weave the product mention into a helpful
    answer if and only if the product is genuinely relevant. Every sponsored
    mention carries a \texttt{[Sponsored]} disclosure tag, ensuring compliance
    with Federal Trade Commission (FTC) endorsement guidelines.

  \item \textbf{Multi-dimensional response and engagement evaluation.} Beyond
    conventional impression and click-through tracking, a dedicated LLM-as-judge
    module analyses each post-ad follow-up message and returns a structured set
    of metrics: an engagement type classification (product inquiry, comparison,
    price inquiry, purchase intent, feature question, positive or negative
    reaction, dismissal), a continuous engagement score, sentiment toward the
    advertisement, a naturalness score capturing how intrusive the placement
    felt, and an estimated purchase-proximity score. These metrics jointly
    form a higher-fidelity signal than any single click-or-not-click
    measurement and are the substrate on which downstream optimisation
    operates.

  \item \textbf{Ethical safeguards.} Conversations classified as sensitive. They include those
    involving health crises, emotional distress, financial hardship, or
    grief and are automatically excluded from ad injection, regardless of
    commercial relevance.

  \item \textbf{Adaptive per-ad context optimisation.} Each advertisement in
    the inventory carries a dynamic presentation-guidance document which is injected into the response-synthesis prompt
    alongside the product description. An initial version of this context is
    generated automatically by an LLM at the moment the ad is created, so that
    every ad begins life with a coherent presentation strategy rather than an
    empty field. A background optimiser then periodically reviews the
    accumulated engagement metrics for each ad, diagnoses what is working and
    what is not (low naturalness, high dismissal rate, weak purchase proximity
    despite strong topical match), and asks a reasoning LLM to rewrite the
    context guide accordingly. Every revision is versioned and persisted to an
    \texttt{ad\_context\_history} table together with the metrics snapshot that
    triggered it and a natural-language justification, yielding a fully
    auditable record of how each ad's presentation strategy evolved over time.

  \item \textbf{Image generation with organic product placement.} The platform
    exposes a second generative surface alongside text chat: a user-prompted
    image-generation endpoint that routes through the same context extraction
    and ad matching pipeline. When a relevant sponsored product is identified,
    the user's original image prompt is rewritten by an auxiliary LLM into a
    placement prompt that fits the product into the requested scene, and the
    final image is rendered by a diffusion model (\texttt{gpt-image-1}). When
    the advertised product has a reference photograph, it is supplied to the
    image model via an edit-mode call so that the rendered product is visually
    grounded in the real article rather than hallucinated from its textual
    description.

  \item \textbf{Advertiser portal with campaign and budget management.} A
    self-serve REST API and accompanying React interface allow advertisers
    to register, organise their spend into campaigns with optional start
    dates, end dates, and daily budget caps, and create ads within each
    campaign. Campaigns follow an explicit lifecycle with
    auto-completion when the total budget is exhausted, and daily spend is
    reset lazily at first touch each new day rather than via a cron job.
    Ad creation is accelerated by a URL-based extraction flow in which the
    advertiser pastes a product page URL, the server scrapes it, and an LLM
    call populates the ad fields (name, description, creative copy, target
    topics and intents, brand-safety tags) as a pre-filled draft that the
    advertiser can review and edit before submission.

  \item \textbf{Dual analytics dashboards.} The platform ships two distinct
    analytics surfaces backed by a shared event log. A global, platform-wide
    dashboard reports aggregate impressions, click-through rate, engagement
    events, estimated revenue, and top-performing ads across all advertisers,
    intended for the platform operator. An advertiser-facing dashboard
    reports the same metrics scoped to a single advertiser or a single
    campaign, with a daily breakdown driven by the \texttt{daily\_spend\_log}
    table and visual indicators of budget utilisation.

  \item \textbf{Functional MVP.} The deliverable is a working prototype that
    demonstrates the complete pipeline from user message to ad-augmented
    response across both text and image modalities. Together with the
    advertiser portal, the adaptive context optimiser, and both analytics
    dashboards.
\end{enumerate}

\subsection{Scope and Constraints}
\label{subsec:scope}

The platform is scoped as a minimum viable product (MVP) and is subject to
several deliberate constraints that bound the design space.

The architecture is a single-service deployment: one FastAPI orchestrator
coordinates all internal modules, a single Redis instance manages session state,
and a single PostgreSQL database (extended with \texttt{pgvector}) stores the
ad inventory, the advertiser and campaign records, the daily spend log, the
ad-context history, and the event log. This monolithic design is appropriate
for the current scale, moderate request
concurrency and avoids the operational complexity of a distributed
microservice architecture.

Two external dependencies impose hard constraints on latency and cost. Every
chat request requires at minimum two LLM API calls (one lightweight call for
context extraction and one full-model call for response synthesis) plus one
embedding API call for the query vector. The total end-to-end latency budget
is three seconds, inclusive of all network round-trips and internal
processing. The per-request cost is dominated by these API calls, and the
re-ranking stage uses heuristic scoring weights (a weighted combination of
semantic similarity, normalised bid price, and purchase signal) rather than a
learned model, since no click-through data is available at launch. Image
generation is treated as a separate interactive surface with its own, looser
latency budget, since diffusion-model calls routinely exceed ten seconds and
are dominated by GPU inference time rather than by the orchestration overhead
that constrains the text path.

The engagement evaluator is implemented as a single LLM call that returns the
structured metrics described in Section~\ref{subsec:objectives}. It runs
asynchronously, off the critical request path, so its latency does not count
against the three-second response budget; the trade-off is that engagement
metrics become available only after a short delay rather than synchronously
with the turn in which the ad was shown.

The adaptive context optimiser is similarly a deferred, out-of-band process.
Optimisation for a given advertisement is gated by two guards. A minimum of
five recorded engagements and a minimum interval of twenty-four hours since
the previous revision, which together prevent the system from overfitting to
noise and from thrashing the ad context under light traffic. An advertisement
with insufficient traffic therefore retains its LLM-generated initial context
indefinitely, and optimisation only begins once the statistical foundation is
adequate.

The advertiser portal is scoped to the operational primitives needed to run a
campaign end-to-end: advertiser registration, campaign CRUD with budget and
schedule controls, ad CRUD with URL-based pre-filling, and scoped analytics.
Authentication middleware, invoicing, payment processing, and advertiser
self-service account recovery are deliberately out of scope for the MVP; the
portal assumes a trusted operator environment and defers commercial-grade
identity and billing to subsequent iterations.

\subsection{Report Organisation}
\label{subsec:report-org}

The remainder of this report is organised as follows.
Section~\ref{sec:literature-review} surveys the relevant literature across four
areas: the evolution of digital advertising formats, the capabilities of large
language models for controlled generation, retrieval-augmented generation as a
paradigm for context-aware information retrieval, and the emerging body of work
on advertising within conversational interfaces. Section~\ref{sec:problem-statement}
formulates the problem precisely, identifying the technical challenges of
real-time intent extraction, low-latency semantic ad matching, natural ad
integration, engagement measurement, and per-campaign budget enforcement,
together with the ethical and regulatory constraints that any solution must
satisfy.
Section~\ref{sec:solution-design} presents the system architecture and the
design of each component: the FastAPI orchestrator, the context extraction
module, the Redis-backed session manager, the \texttt{pgvector}-based ad
matching pipeline with budget-gated eligibility, the prompt-engineered
response synthesis module, the image-generation endpoint with
product-placement prompt rewriting, the LLM-as-judge engagement evaluator,
the adaptive ad-context optimiser, the advertiser portal with its
URL-based ad extraction flow, and the dual analytics subsystem serving both
the platform-wide and advertiser-scoped dashboards. 
Section~\ref{sec:illustrative-examples} walks through end-to-end examples
that exercise representative paths through the pipeline, including successful
ad injection, sensitive-conversation suppression, irrelevant-ad rejection,
engagement follow-up detection, multi-turn context accumulation, an
image-generation request with organic product placement, a campaign reaching
its daily budget cap, and an advertiser creating a new ad via URL extraction.
Section~\ref{sec:conclusions} summarises the contributions, discusses
limitations, and outlines directions for future work.